% =====================================================================
%  Centralized Multi-Period OPF (COPF): non-convex, MPISOCP, linear
%  Companion to dddp_otd_section.tex.
%  Requires: amsmath, amssymb, booktabs, float, algorithm, algpseudocode, bm
% =====================================================================

\section{Centralized Multi-Period Optimal Power Flow}
\label{sec:copf}

We study the day-ahead operation of an unbalanced radial distribution
network (DN) that hosts distributed energy resources (DERs) and battery
energy storage systems (BESS). Over a scheduling horizon
$\mathcal{T}=\{1,\dots,T\}$, the operator seeks the DER, storage, and
power-flow schedule that minimizes the cost of energy imported at the
substation while honoring the network physics and all device limits.
Because each battery couples consecutive timesteps through its
state of charge (SOC), the problem is a single, temporally coupled
\emph{multi-period} OPF rather than a sequence of independent snapshots.

This section develops three models that will later serve as the atomic
building block inside each temporal window of the decomposition in
Section~\ref{sec:dddp}. Section~\ref{ssec:notation} fixes the notation
and network model; Section~\ref{ssec:nonconvex} states the exact
non-convex formulation; Section~\ref{ssec:mpisocp} presents the
Multi-Period Iterative SOCP (MPISOCP) scheme that solves it to global
accuracy; and Section~\ref{ssec:linear} gives the loss-free linear
counterpart used when speed is favored over exactness.

% ---------------------------------------------------------------
\subsection{Notation and Network Model}
\label{ssec:notation}

The DN is modeled as a directed graph $G=(\mathcal{N},\mathcal{E})$,
where $\mathcal{N}$ and $\mathcal{E}$ denote its nodes and branches. The
superscript $p$ indexes the phase, $\Phi_j$ and $\Phi_{ij}$ collect the
phases present at bus $j$ and branch $(i,j)$, and $t$ indexes time.
For branch $(i,j)$ and phase $p$ at time $t$, $P^{pp}_{ij,t}$ and
$Q^{pp}_{ij,t}$ are the active and reactive power flows. At node $j$,
$p^{p}_{L,j,t}$ and $q^{p}_{L,j,t}$ are the load demands, while
$p^{p}_{D,j,t}$ and $q^{p}_{D,j,t}$ are the DER active and reactive
injections, bounded by the inverter rating $S^{p}_{\mathrm{DR},j}$.
The BESS at node $j$ charges and discharges with powers $P^{p}_{c,j,t}$
and $P^{p}_{d,j,t}$, has stored energy (SOC) $B^{p}_{j,t}$, rated power
$P_{\mathrm{BR},j}$, charge/discharge efficiencies $\eta_c,\eta_d$, and
SOC bounds $[\mathrm{soc}_{\min},\mathrm{soc}_{\max}]$. We work with the
squared voltage magnitude $\upsilon^{p}_{j,t}=|V^{p}_{j,t}|^2$ and, for
the branch-flow model, with the squared current magnitude
$l^{pp}_{ij,t}=|I^{pp}_{ij,t}|^2$ and the cross-phase current product
$l^{pq}_{ij,t}=|I^{pp}_{ij,t}|\,|I^{qq}_{ij,t}|$; $\delta^{pq}_{ij,t}$
denotes the corresponding current-angle difference. Line impedance and
resistance are $z^{pq}_{ij}$ and $R^{pq}_{ij}$. Voltage magnitudes are
constrained to $[V_{\min},V_{\max}]$. The principal symbols are
collected in Table~\ref{tab:notation}.

\begin{table}[H]
  \centering
  \caption{Key Notation}
  \label{tab:notation}
  \begin{tabular}{ll}
    \toprule
    Symbol & Description \\
    \midrule
    \(G=(\mathcal{N},\mathcal{E})\) & Directed graph of nodes and branches \\
    \(\Phi_j,\,\Phi_{ij}\) & Set of phases at bus \(j\) and branch \((i,j)\) \\
    \(P^{pp}_{ij,t},\,Q^{pp}_{ij,t}\) & Active/reactive flow on \((i,j)\), phase \(p\), time \(t\) \\
    \(p^{p}_{L,j,t},\,q^{p}_{L,j,t}\) & Active/reactive demand at node \(j\), phase \(p\) \\
    \(p^{p}_{D,j,t},\,q^{p}_{D,j,t}\) & DER active/reactive injection at node \(j\) \\
    \(S^{p}_{\mathrm{DR},j}\) & DER inverter apparent-power rating \\
    \(P^{p}_{c,j,t},\,P^{p}_{d,j,t}\) & BESS charge/discharge active power \\
    \(B^{p}_{j,t}\) & BESS state of charge (SOC) \\
    \(P_{\mathrm{BR},j}\) & BESS rated active-power capacity \\
    \(\eta_c,\,\eta_d\) & BESS charge/discharge efficiencies \\
    \(\mathrm{soc}_{\min},\,\mathrm{soc}_{\max}\) & BESS SOC bounds \\
    \(V^{p}_{j,t}\) & Complex voltage \((|V^{p}_{j,t}|\angle\theta^{p}_{j,t})\) \\[2pt]
    \(\upsilon^{p}_{j,t}\) & Squared voltage magnitude \((\upsilon^{p}_{j,t}=|V^{p}_{j,t}|^2)\) \\[2pt]
    \(I^{pp}_{ij,t}\) & Complex current \((|I^{pp}_{ij,t}|\angle\delta^{pp}_{ij,t})\) \\[2pt]
    \(\delta^{pq}_{ij,t}\) & Current-angle difference \((\delta^{pp}_{ij,t}-\delta^{qq}_{ij,t})\) \\[2pt]
    \(l^{pp}_{ij,t}\) & Squared current magnitude \((|I^{pp}_{ij,t}|^2)\) \\[2pt]
    \(l^{pq}_{ij,t}\) & Current-magnitude product \((|I^{pp}_{ij,t}|\,|I^{qq}_{ij,t}|)\) \\[2pt]
    \(z^{pq}_{ij},\,R^{pq}_{ij}\) & Impedance/resistance of line \((i,j)\), \(pq\in\Phi_{ij}\) \\
    \bottomrule
  \end{tabular}
\end{table}

% ---------------------------------------------------------------
\subsection{Non-Convex Multi-Period Formulation}
\label{ssec:nonconvex}

The centralized multi-period three-phase OPF minimizes the total
substation energy cost over the horizon,
\begin{equation}\label{eq0}
\small
    \min_{\mathcal{Y}}\; F(\mathcal{Y}),
\end{equation}
where $\mathcal{Y}$ is the set of decision variables and
\begin{equation}\label{eq1}
\begin{small}
\begin{aligned}
F(\mathcal{Y})=\sum_{t=1}^{T} \Bigg[
& C_{t}\,P_{subs,t} \\
& + \alpha \sum_{j \in \mathcal{B}} \left\{ (1 - \eta_c) P_{C_j,t}
+ \left( \tfrac{1}{\eta_d} - 1 \right) P_{D_j,t} \right\}
\Bigg].
\end{aligned}
\end{small}
\end{equation}
The first term is the cost of substation power, with $C_t$
($\$/\mathrm{kWh}$) the price and $P_{subs,t}$ (kW) the total imported
power at step $t$. The second term penalizes BESS conversion losses
with weight $\alpha>0$; as shown in~\cite{Nazir}, this penalty
suppresses simultaneous charging and discharging without introducing
binary variables, keeping the model continuous.

The schedule must satisfy the branch-flow physics of the unbalanced DN.
Active and reactive nodal balance at bus $j$, phase $p$, time $t$ are
\begin{equation}\label{eq2}
\begin{split}
P_{ij,t}^{pp}
= \sum_{(j,k)\in E} P_{jk,t}^{pp} + p_{L_j,t}^p
- p_{D_j,t}^p - P_{d_j,t}^p + P_{c_j,t}^p \\ + \sum_{q \in \Phi_j} l_{ij,t}^{pq}
\left( r_{ij}^{pq} \cos(\delta_{ij,t}^{pq}) + x_{ij}^{pq} \sin(\delta_{ij,t}^{pq}) \right),
\end{split}
\end{equation}
\begin{equation}\label{eq3}
\begin{split}
Q_{ij,t}^{pp}
= & \sum_{(j,k)\in E} Q_{jk,t}^{pp} + q_{L_j,t}^p - q_{D_j,t}^p \\
& + \sum_{q \in \Phi_j} l_{ij,t}^{pq}
\left( x_{ij}^{pq} \cos(\delta_{ij,t}^{pq}) - r_{ij}^{pq} \sin(\delta_{ij,t}^{pq}) \right),
\end{split}
\end{equation}
in which the last summation accounts for the series line losses. The
voltage drop along branch $(i,j)$ follows the three-phase Kirchhoff
voltage law
\begin{equation}\label{eq4}
\begin{split}
v_{j,t}^p =  & v_{i,t}^p
- \sum_{q \in \Phi_j} 2\,\mathbb{R}\!\left[ S_{ij,t}^{pq} \left( z_{ij}^{pq} \right)^* \right]- \\ & \sum_{\substack{q_1, q_2 \in \Phi_j}}
2\,\mathbb{R}\!\left[ z_{ij}^{pq_1} l_{ij,t}^{q_1 q_2}
\left( \angle(\delta_{ij,t}^{q_1 q_2}) \right)
\left( z_{ij}^{pq_2} \right)^* \right].
\end{split}
\end{equation}
The exactness of the branch-flow representation is enforced by the
power--voltage--current relation and the cross-phase current coupling,
\begin{equation}\label{eq5}
\left( P_{ij,t}^{pp} \right)^2 + \left( Q_{ij,t}^{pp} \right)^2
= v_{i,t}^p \, l_{ij,t}^{pp},
\end{equation}
\begin{equation}\label{eq6}
\left( l_{ij,t}^{pq} \right)^2 = l_{ij,t}^{pp} \, l_{ij,t}^{qq}.
\end{equation}

Device and operating limits complete the model. DER reactive capability
is bounded by the inverter rating,
\begin{equation}\label{linqder}
{ - \sqrt {{{\left( {S_{D\,R_j}^p} \right)}^2} - {{\left( {p_{D_j,t}^p} \right)}^2}} \le q_{D_j,t}^p \le \sqrt {{{\left( {S_{D\,R_j}^p} \right)}^2} - {{\left( {p_{D_j,t}^p} \right)}^2}}},
\end{equation}
and bus voltages are kept within statutory bounds,
\begin{equation}\label{linv}
{\left( {{V_{min}}} \right)^2} \le {\upsilon}_j^p \le {\left( {{V_{max}}} \right)^2}.
\end{equation}
The BESS is governed by its SOC dynamics, power limits, energy bounds,
and a cyclic terminal condition,
\begin{equation}\label{linsocevol}
B_{j,t}^{p} = B_{j,t-1}^{p} + \Delta t\, \eta_c P_{c_j,t}^{p} - \Delta t\, (1/\eta_d) P_{d_j,t}^{p},
\end{equation}
\begin{equation}\label{linPch}
P_{cj,t}^{p} \le P_{BR_j},
\end{equation}
\begin{equation}\label{linPdis}
P_{dj,t}^{p} \le P_{BR_j},
\end{equation}
\begin{equation} \label{linsoclim}
soc_{min}B_{R,j} \le B_{j,t}^p \le  soc_{max}B_{R,j},
\end{equation}
\begin{equation}\label{lintermsoc}
B_{j,0}^{p} = B_{j,T}^{p}.
\end{equation}

The SOC recursion~\eqref{linsocevol} is what couples the timesteps into
a single multi-period program: every $B^{p}_{j,t}$ with $t\ge 1$ is a
decision variable, while at the first step $t=1$ the preceding state is
fixed to the given initial SOC, $B^{p}_{j,t-1}=B^{p}_{j,0}$, which
anchors the entire horizon. Constraints~\eqref{linPch}--\eqref{linPdis}
cap the charge/discharge power, \eqref{linsoclim} keeps the stored
energy within its usable band, and~\eqref{lintermsoc} restores the SOC
to its initial value at the end of the horizon so that scheduling one
day does not deplete the batteries for the next.

\paragraph*{Source of non-convexity and the current-angle
simplification.} Treating the current angles $\delta^{pq}_{ij,t}$ as
variables renders~\eqref{eq2}--\eqref{eq4} strongly nonlinear and
non-convex. In distribution networks, however, the current angles
deviate only marginally from their base-case values~\cite{XX}. We
therefore fix $\delta^{pq}_{ij,t}$ to the angles obtained from a
base-case OpenDSS power flow, held constant across all timesteps. With
this simplification, \eqref{eq2}--\eqref{eq4} become linear in the
remaining variables. The only residual non-convexities are the equality
constraints~\eqref{eq5}--\eqref{eq6}, which remain non-convex because of
their bilinear right-hand sides; these are handled next.

% ---------------------------------------------------------------
\subsection{MPISOCP: Iterative SOCP Convexification}
\label{ssec:mpisocp}

We convexify the remaining non-convex equalities with a Multi-Period
Iterative Second-Order Cone Programming (MPISOCP) scheme. The
equalities~\eqref{eq5}--\eqref{eq6} are first relaxed into second-order
cones,
\begin{equation}\label{eqpvc}
\left( P_{ij,t}^{pp} \right)^2 + \left( Q_{ij,t}^{pp} \right)^2
\leq v_{i,t}^p \, l_{ij,t}^{pp},
\end{equation}
\begin{equation}\label{eqcmc}
\left( l_{ij,t}^{pq} \right)^2 \leq l_{ij,t}^{pp} \, l_{ij,t}^{qq}.
\end{equation}
The inequalities enlarge the feasible set and make each stage a convex
conic program. The relaxation is exact only when both cones are tight,
i.e.\ when the gaps between the two sides of~\eqref{eqpvc}
and~\eqref{eqcmc} vanish. We quantify these gaps at every branch,
phase, and time by the residuals
\begin{equation}\label{erroreqpvc}
\left( E1_{ij,t}^{pp} \right) = \left( P_{ij,t}^{pp} \right)^2 + \left( Q_{ij,t}^{pp} \right)^2
- v_{i,t}^p \, l_{ij,t}^{pp},
\end{equation}
\begin{equation}\label{erroreqcmc}
\left( E2_{ij,t}^{pq} \right) = \left( l_{ij,t}^{pq} \right)^2 - l_{ij,t}^{pp} \, l_{ij,t}^{qq},
\end{equation}
both of which are non-positive by construction. MPISOCP drives them to
zero over successive iterations. Because the residuals are non-positive,
tightening at iteration $m{+}1$ is imposed by requiring
\begin{equation}\label{erroreqpvcred}
(E1_{ij,t}^{pp})^{m+1} \geq \tau\, (E1_{ij,t}^{pp})^{m},
\end{equation}
\begin{equation}\label{erroreqcmcred}
(E2_{ij,t}^{pq})^{m+1} \geq \tau\, (E2_{ij,t}^{pq})^{m},
\end{equation}
with $\tau\in[0,1]$: each residual is pushed to at least a fraction
$\tau$ of its previous (negative) value, moving it toward zero. The
right-hand sides of~\eqref{erroreqpvcred}--\eqref{erroreqcmcred} are
constants known after iteration $m$, but the left-hand sides are still
bilinear. Linearizing them by a first-order Taylor expansion about the
iterate-$m$ operating point yields the tractable directional cuts
\begin{equation}\label{eqdirpvc}
\begin{split}
2P_{ij,t}^{pp(m)} P_{ij,t}^{pp(m+1)}
&+ 2Q_{ij,t}^{pp(m)} Q_{ij,t}^{pp(m+1)}
- l_{ij,t}^{pp(m)} v_{i,t}^{p(m+1)} \\&
- v_{i,t}^{p(m)} l_{ij,t}^{pp(m+1)}
\geq (\tau + 1)\, E1_{ij,t}^{pp(m)},
\end{split}
\end{equation}
\begin{equation}\label{eqdircmc}
\begin{split}
2l_{ij,t}^{pq(m)}  l_{ij,t}^{pq(m+1)}
- l_{ij,t}^{pp(m)}  l_{ij,t}^{qq(m+1)}
- l_{ij,t}^{qq(m)}  l_{ij,t}^{pp(m+1)} \\
\geq (\tau + 1)\, E2_{ij,t}^{pq(m)}.
\end{split}
\end{equation}
These linear constraints steer the conic solution toward the exact
manifold, shrinking the relaxation gap monotonically. The complete
procedure is summarized in Algorithm~\ref{alg:isocp}: it repeatedly
solves the SOCP relaxation augmented with the current directional cuts,
re-linearizes at the new iterate, and terminates once the maximum
residual falls below a tolerance $\varepsilon$ or an iteration cap
$M_{\max}$ is reached.

\begin{algorithm}[t]
\caption{Multi-Period Iterative SOCP (MPISOCP).}
\label{alg:isocp}
\begin{algorithmic}[1]
\Require $\tau\!\in\![0,1]$, tolerance $\varepsilon$, cap $M_{\max}$.
\State Solve the SOCP relaxation
\eqref{eq0}--\eqref{eq4},\,\eqref{eqpvc},\,\eqref{eqcmc}, \eqref{linqder}--\eqref{lintermsoc}; obtain
$\bm{z}^{(0)}\!=\!(P^{(0)},Q^{(0)},v^{(0)},l^{(0)})$.
\State Compute $E1^{pp(0)}_{ij,t}$, $E2^{pq(0)}_{ij,t}$
from~\eqref{erroreqpvc}--\eqref{erroreqcmc} and
$E^{(0)}\!\leftarrow\!\max\{|E1^{pp(0)}|,|E2^{pq(0)}|\}$.
\For{$m=0,1,\dots,M_{\max}\!-\!1$}
    \If{$E^{(m)}\!\le\!\varepsilon$} \Return $\bm{z}^{(m)}$. \EndIf
        \State Linearize the residuals at $\bm{z}^{(m)}$ to form the
directional cuts~\eqref{eqdirpvc}--\eqref{eqdircmc}.
\State Solve \eqref{eq0}--\eqref{eq4},\,\eqref{eqpvc},\,\eqref{eqcmc},\eqref{linqder}--\eqref{lintermsoc}, augmented with the cuts~\eqref{eqdirpvc}--\eqref{eqdircmc}; obtain
$\bm{z}^{(m+1)}$ and update $E^{(m+1)}$.
\EndFor
\State \Return $\bm{z}^{(M_{\max})}$.
\end{algorithmic}
\end{algorithm}

% ---------------------------------------------------------------
\subsection{Linear Multi-Period Formulation}
\label{ssec:linear}

When solution speed is prioritized over loss accuracy, the quadratic
loss terms in the balance and voltage equations are dropped, yielding a
fully linear multi-period OPF. It retains the objective~\eqref{eq1} and
all device constraints, but replaces~\eqref{eq2}--\eqref{eq4} with their
loss-free counterparts:
\begin{equation}\label{linPij}
\min_{\mathcal{Y}} F(\mathcal{Y}) \;\text{ s.t. }\; P_{ij,t}^{pp} = \sum_{(j,k)\in E} {P_{jk,t}^{pp}} + p_{{L_j,t}}^p - p_{D_j,t}^p - P_{d_j,t}^p + P_{c_j,t}^p,
\end{equation}
\begin{equation}\label{linQij}
Q_{ij,t}^{pp} = \sum_{(j,k)\in E} {Q_{jk,t}^{pp}} + q_{{L_j,t}}^p - q_{Dj,t}^p,
\end{equation}
\begin{equation}\label{linkvl}
{\upsilon }_{j,t}^p = {\upsilon}_{i,t}^p - \sum\limits_{q \in {\phi _j}} {2\left[ {S_{ij,t}^{pq}\left( {z_{ij}^{pq}} \right)^*} \right]},
\end{equation}
together with the DER and BESS
constraints~\eqref{linqder}--\eqref{lintermsoc}. This linear model has
no bilinear terms and hence needs no convexification; it is solved
directly as a single linear program and serves as the fast alternative
to MPISOCP inside the temporal decomposition of Section~\ref{sec:dddp}.

% =====================================================================
%  End of section
% =====================================================================
